\documentclass{ltjsbook}
\usepackage{docmute} 
\usepackage{../../myStyle} 

\begin{document}
\chapter{ベクトルの空間的理解}
\label{L04_Space}
\section{空間としてのベクトル}
ここまで代数(文字と式)としての行列の話をしてきましたが，ここからは空間的なイメージで考えていきたいと思います。
何かを理解するときの入り口はいろいろあった方がよいと思います。数字の計算では面白みがわからなかった人も，空間イメージであればわかりやすいということがあるかもしれません。もちろん数字の羅列が大好きな人はそのままでもいいのですが，空間イメージと合わせて考えることができるという，数学世界の完璧なまでの整合性に感動するのも一興ではないでしょうか。

心理統計学的な観点からみると，線形代数は多変量解析に用いるツールということになります。
しかし線形代数は，たとえばコンピュータグラフィクスの世界でも必要な知識です。CGというと空間をさまざまな角度から描画するような，視点の移動や空間の回転などが特徴的かとおもいますが，この空間の回転は行列の仕事です。
また，物理学では物質に働く力，その大きさと向きを考えますが，これはベクトルとして表現されます。
線形代数の数学的側面はこのように，統計学や情報科学，物理学の間にあってそれぞれの世界を媒介する共通言語なのです。

それではさっそく，ベクトルの空間的イメージをつかむところから始めましょう。
\subsection{ベクトル空間}

\subsection{}

ベクトルの和・差
積（スケイラー）
ベクトルの長さ

% 4. 空間的理解
% - ベクトルの内積
% - 内積と相関
% - ベクトルの直交

\end{document}